\section{Integral Curves and Flows}

\subsection{Integral curves}

\bdefn

    Let $V$ be a vector field on a smooth manifold $M$.
    An {\emphcolor integral curve} for $V$ is a smooth curve $\gamma\colon J\to M$ (with $J$ an interval in the real line)
    such that $\gamma'(t)=V_{\gamma(t)}$.
    If $0\in J$, then $\gamma(0)$ is called the {\emphcolor starting point} of $\gamma$.

\edefn

Recall that $\gamma'(t)\in T_{\gamma(t)}M$ is viewed as the derivation
$$ \gamma'(t)f = (f\circ\gamma)'(t) . $$
Equivalently, the differential $d\gamma_t$ is a map $\bR\cong T_tJ\to T_{\gamma(t)}M$, and then $\gamma'(t)f=d\gamma_t(1)(f)$, where $1\in\bR$
corresponds to the directional derivative in the direction of $1$ (the usual $1$-dimensional derivative).

\bnote

    In this section we will use the notation $\dot\gamma(t)$ for the derivative of a curve with respect to $t$.
    This is because often, the superscript is already being used as an index.

\enote

\bexam

    Consider $\bR^2$ and let $V=\ipdv x$ be the first-coordinate vector field.
    A curve $\gamma(t)=(\gamma^1(t),\gamma^2(t))$ has derivative $\gamma'(t)=\dot\gamma^1(t)\ipdv x+\dot\gamma^2(t)\ipdv y$.
    Then $\gamma$ is an integral curve of $V$ when it is of the form $\gamma(t)=(a+t,b)$.

\eexam

Finding integral curves boils down to solving ODEs.
For a vector field $V$, let $\gamma\colon J\to M$ be a smooth curve.
On a coordinate domain $U\subseteq M$ we can write $\gamma$ in local coordinates $\gamma(t)=(\gamma^1(t),\dots,\gamma^n(t))$ as well as
$V$ as $V^i(p)\evpdv{x^i}_p$.
Then the requirement that $\gamma'(t)=V_{\gamma(t)}$ is just
$$ \dot\gamma^i(t)\evpdv{x^i}_{\gamma(t)} = V^i(\gamma(t))\evpdv{x^i}_{\gamma(t)} , $$
which is just the system of ODEs:
$$ \eqalign{
    \dot\gamma^1(t) &= V^1(\gamma^1(t),\dots,\gamma^n(t))\cr
    &\vdots\cr
    \dot\gamma^n(t) &= V^1(\gamma^1(t),\dots,\gamma^n(t))\cr
} $$

Recall the existence and uniqueness of a solution to ODEs.
Since everything here is smooth, so we can derive the following simple consequence.

\bprop

    Let $V$ be a smooth vector field on $M$.
    Then for every $p\in M$ there exists an $\epsilon>0$ and a unique smooth curve $\gamma\colon(-\epsilon,\epsilon)\to M$ that is an integral
    curve of $V$ starting at $p$.

\eprop

\blemm

    Let $V$ be a smooth vector field on $M$, $J\subseteq\bR$ an interval, and $\gamma\colon J\to\bR$ an integral curve of $V$.
    For any $a\in\bR$, the curve $\tilde\gamma\colon a^{-1}J\to M$ defined by $\tilde\gamma(t)=\gamma(at)$ is an integral curve of $aV$.

    Similarly $\hat\gamma\colon J-a\to M$ defined by $\hat\gamma(t)=\gamma(t+a)$ is an integral curve of $V$.

\elemm

\bproof

    We look at how $\tilde\gamma$ acts on $C^\infty$ functions defined in a neighborhood of $t_0$.
    $$ \tilde\gamma'(t_0)f = \evdrv t_{t_0}(f\circ\tilde\gamma)(t) = \evdrv t_{t_0}(f\circ\gamma)(at) , $$
    which by the chain rule is just
    $$ = a(f\circ\gamma)'(at_0) = a\gamma'(at_0)f = aV_{\gamma(at_0)}f = aV_{\tilde\gamma(t_0)}f . $$
    The translation $\hat\gamma$ is proven similarly.

\eproof

\bprop

    If $F\colon M\to N$ is smooth, then $X\in\X(M),Y\in\X(N)$ are $F$-related iff $F$ takes integral curves of $X$ to integral curves of $Y$.
    Meaning if $\gamma$ is integral for $X$, then $F\circ\gamma$ is integral for $Y$.

\eprop

\bproof

    If $X$ and $Y$ are $F$-related, let $\gamma\colon J\to M$ be integral for $X$, and define $\sigma\colon J\to N$ by $\sigma=F\circ\gamma$.
    Then
    $$ \sigma'(t) = (F\circ\gamma)'(t) = dF_{\gamma(t)}(\gamma'(t)) = dF_{\gamma(t)}(V_{\gamma(t)}) = Y_{F\circ\gamma(t)} = Y_{\sigma(t)}, $$
    as desired.

    Conversely suppose $F$ maps integral curves of $X$ to integral curves of $Y$.
    Let $p\in M$ and let $\gamma\colon J\to M$ be an integral curve of $X$ starting at $p$.
    Then $\sigma=F\circ\gamma$ is integral for $Y$ so
    $$ Y_{F(p)} = (F\circ\gamma)'(0) = dF_{\gamma(0)}(\gamma'(0)) = dF_p(X_p), $$
    meaning $X$ and $Y$ are $F$-related.
    \qed

\eproof

\subsection{Flows}

Let $M$ be a smooth manifold and $V\in\X(M)$.
Now suppose that for every $p\in M$, $V$ has a unique integral curve starting at $p$ which is defined for all $t\in\bR$ (this is a big assumption,
but will suffice for the explanation), denoted by $\theta^{(p)}\colon\bR\to M$.
Now for each $t\in\bR$ we define $\theta_t\colon M\to M$ by mapping $p\in M$ to $\theta^{(p)}(t)$.
In other words we have
$$ \theta\colon\bR\times M\to M . $$
Translation implies that $t\mapsto\theta^{(p)}(s+t)$ is an integral curve for $V$ starting at $q=\theta^{(p)}(s)$.
Uniqueness implies that $\theta^{(q)}(t)=\theta^{(p)}(t+s)$, so that
$$ \theta_t\circ\theta_s(p) = \theta_t(\theta^{(p)}(s)) = \theta^{(q)}(t) = \theta_{s+t}(p) . $$
Since $\theta_0(p)=\theta^{(p)}(0)=p$, we see that $\theta\colon\bR\times M\to M$ is an action of the additive group $\bR$ on $M$.

This motivates the following definition.

\bdefn

    A {\emphcolor global flow} on $M$ is a continuous left action of $\bR$ on $M$.
    In other words, a continuous map
    $$ \theta\colon\bR\times M\to M $$
    satisfying
    $$ \theta(0,p) = p,\qquad \theta(t+s,p) = \theta(t,\theta(s,p)) . $$

\edefn

\bnote

    Global flows are also called {\emphcolor one-parameter group actions}.

\enote

Given a global flow $\theta$ on $M$ we define
$$ \theta_t\colon M\to M,\ \theta_t(p)=\theta(t,p),\qquad \theta^{(p)}\colon\bR\to M,\ \theta^{(p)}(t)=\theta(t,p) , $$
as with Lie group actions.
The image of the curve $\theta^{(p)}$ is the orbit of $p$ under the action $\theta$.
And the requirement that $\theta$ be an action is equivalent to requiring
$$ \theta_{t+s} = \theta_t\circ\theta_s . $$

If $\theta\colon\bR\times M\to M$ is a global flow, define a tangent vector $V_p\in T_pM$ by
$$ V_p = {\theta^{(p)}}'(0) . $$
The assignment $p\mapsto V_p$ is a rough vector field on $M$, called the {\emphcolor infinitesimal generator} of $\theta$.

\bprop

    Let $\theta$ be a smooth global flow on $M$.
    Its infinitesimal generator $V$ is then a smooth vector field on $M$, and each $\theta^{(p)}$ is an integral curve of $V$ starting at $p$.

\eprop

\bproof

    To show that $V$ is smooth, it is sufficient to show that $Vf$ is smooth for all $f\in C^\infty(M)$.
    Now, for $p\in M$,
    $$ Vf(p) = V_pf = {\theta^{(p)}}'(0)f = \evdrv t_0f(\theta^{(p)}(t)) = \evpdv t_{(0,p)}f(\theta(t,p)) . $$
    Since $f\circ\theta\colon\bR\times M\to\bR$ is smooth by assumption, its partial derivatives are smooth as well.
    Thus $Vf$ is smooth, as desired.

    Now we need to show that $\theta^{(p)}$ is an integral curve of $V$.
    For $t\in\bR$, let $q=\theta^{(p)}(t)=\theta(t,p)$ so that
    $$ V_{\theta^{(p)}(t)} = V_q = {\theta^{(q)}}'(0) . $$
    We want to show that this is equal to ${\theta^{(p)}}'(t)$.
    Now,
    $$ \theta^{(q)}(s) = \theta(s,q) = \theta(s,\theta(t,p)) = \theta(t+s,p) = \theta^{(p)}(t+s) . $$
    And so we have that
    $$ {\theta^{(q)}}'(0) = {\theta^{(p)}}'(t), $$
    as desired.
    \qed

\eproof

While every global flow gives rise to a vector field, its infinitesimal generator, the converse is not true in general.
The issue is the globalness: a vector field's integral curve does not in general need to be defined at every point on the real line.

\bexam

    Let $M=\bR^2-0$, and $V=\ipdv x$.
    The unique integral curve starting at $(-1,0)$ is $\gamma(t)=(-1+t,0)$.
    But $\gamma$ then cannot be defined at $t=1$, as this is a removable singularity and defining it at $t=1$ would require defining it to be equal to $0\notin M$.

    Another example is $M=\bR^2$ and $V=x^2\ipdv x$.
    An integral curve $\gamma(t)=(\gamma^1(t),b)$ must satisfy $\dot\gamma^1(t)=\gamma^1(t)^2$ and thus has the solution $\gamma^1(t)=1/(a-t)$.
    If we require it to start at $(1,0)$, then we have
    $$ \gamma(t) = \parens{\frac1{1-t},0} . $$
    This also cannot be defined at $t=1$.

\eexam

Thus we must somehow localize the idea of a global flow.

\bdefn

    Let $M$ be a smooth manifold.
    A {\emphcolor flow domain} for $M$ is an open subset $\D\subseteq\bR\times M$ with the property that for every $p\in M$, the set $\D^{(p)}=\set{t\in\bR}[(t,p)\in\D]$
    is an open interval containing $0$.

    A {\emphcolor (local) flow} on $M$ is a continuous map $\theta\colon\D\to M$ where $\D\subseteq\bR\times M$ is a flow domain, which satisfies the usual group action laws
    when defined.
    Explicitly, for all $p\in M$, $\theta(0,p)=p$ and
    $$ \theta(t,\theta(s,p)) = \theta(t+s,p) , $$
    for all $s\in\D^{(p)}$, and $t\in\D^{(\theta(s,p))}$ such that $s+t\in\D^{(p)}$.

    As with global flows, we define $\theta_t(p)=\theta^{(p)}(t)=\theta(t,p)$ when $(t,p)\in\D$.
    For $t\in\bR$ we define
    $$ M_t = \set{p\in M}[(t,p)\in\D] $$
    so that $p\in M_t$ iff $t\in\D^{(p)}$ iff $(t,p)\in\D$, and $\theta^{(p)}\colon\D^{(p)}\to M$, $\theta_t\colon M_t\to M$.

    If $\theta$ is smooth, its {\emphcolor infinitesimal generator} is defined by $V_p={\theta^{(p)}}'(0)$.

\edefn

\bnote

    The term {\emphcolor local one-parameter group action} is also used for local flows.

\enote

\bprop

    The infinitesimal generator of a flow is a smooth vector field, and $\theta^{(p)}$ is an integral curve for $V$ for each $p\in M$.

\eprop

\bproof

    The proof is the same as in the global case.
    Just verify that all terms are defined locally (in close enough neighborhoods).
    \qed

\eproof

\bdefn

    A {\emphcolor maximal integral curve} is an integral curve which cannot be extended to a larger interval.
    A {\emphcolor maximal flow} is a flow that cannot be extended to a larger flow domain.

\edefn

\bthrm[title=Fundamental theorem on flows]

    Let $V$ be a smooth vector field on $M$.
    Then there is a unique maximal flow $\theta\colon\D\to M$ whose infinitesimal generator is $V$.
    The flow has the following properties:
    \benum
        \item For each $p\in M$, the curve $\theta^{(p)}\colon\D^{(p)}\to M$ is the unique maximal integral curve of $V$ starting at $p$.
        \item If $s\in\D^{(p)}$ then $\D^{(\theta(s,p))}$ is the interval $\D^{(p)}-s=\set{t-s}[t\in\D^{(p)}]$.
        \item For each $t\in\bR$, $M_t$ is open in $M$ and $\theta_t\colon M_t\to M_{-t}$ is a diffeomorphism with inverse $\theta_{-t}$.
    \eenum

\ethrm

\bproof

    Suppose that $\gamma,\tilde\gamma\colon J\to M$ are two integral curves of $V$ defined on the same interval such that they agree at some point $t_0\in J$.
    Then let $\c S=\set{t\in J}[\gamma(t)=\tilde\gamma(t)]$.
    This is nonempty since $t_0\in\c S$ by assumption, and it is closed by continuity (for $t\notin\c S$ then by Hausdorff there are open sets around
    $\gamma(t),\tilde\gamma(t)$ which don't intersect, their preimages under $\gamma,\tilde\gamma$ are open and the intersection doesn't intersect $\c S$).
    But $\c S$ is also open: for $t_1\in\c S$, $\gamma,\tilde\gamma$ are both solutions to the same ODE with the initial condition $\gamma(t_1)=\tilde\gamma(t_1)$,
    and thus agree in a neighborhood of $t_1$.
    Since $J$ is connected, $\c S$ must be equal to $J$.

    Thus any two integral curves defined on the same interval must be equal.
    So let us define $\D^{(p)}$ to be the union of all open intervals $J$ with $0\in J$ such that there is an integral curve of $V$ starting at $p$ defined on $J$.
    Then for $t\in\D^{(p)}$, define $\theta^{(p)}(t)=\gamma(t)$ where $\gamma$ is an integral curve of $V$ starting at $p$ defined at $t$.
    This is well-defined by our above argument, and clearly the unique maximal integral curve starting at $p$.

    Now let $\D=\set{(t,p)\in\bR\times M}[t\in\D^{(p)}]$, and define $\theta\colon\D\to M$ by $\theta(t,p)=\theta^{(p)}(t)$.
    $\theta$ satisfies the first point in our theorem statement.
    Now we need to verify that it satisfies the group laws.
    Clearly $\theta(0,p)=p$.
    For $p\in M$ and $s\in\D^{(p)}$ let $q=\theta(s,p)$, then the curve $\gamma\colon\D^{(p)}-s\to M$ defined by $\gamma(t)=\theta^{(p)}(t+s)$ starts at $q$
    and is also integral (by translation).
    By uniqueness, $\gamma$ and $\theta^{(q)}$ agree on the intersection of their domains, which is equivalent to the second group law:
    $$ \theta^{(p)}(t+s) = \theta^{(q)}(t) \iff \theta(t+s,p) = \theta(t,\theta(s,p)) . $$

    The maximality of $\theta^{(q)}$ further implies that $\gamma$'s domain is contained in $\D^{(q)}$.
    That is, $\D^{(p)}-s\subseteq\D^{(q)}$.
    Since $0\in\D^{(p)}$ this implies $-s\in\D^{(q)}$ and $\theta^{(q)}(-s)=\theta(-s,\theta(s,p))=p$.
    So if we repeat the above argument with $p=\theta(-s,q)$, we see that $\D^{(q)}+s\subseteq\D^{(p)}$ so that $\D^{(q)}-s=\D^{(p)}$, proving the second point.

    Now we show that $\D\subseteq\bR\times M$ is open, so that it is a flow domain, and that $\theta\colon\bR\times M\to M$ is smooth.
    Define $W\subseteq\D$ to be the set of all $(t,p)\in\D$ where $\theta$ is smooth on a product neighborhood of $(t,p)$ of the form $J\times U\subseteq\D$
    where $0\in J$ is a neighborhood and $U\subseteq M$ is a neighborhood of $p$.
    Then $W$ is an open subset of $\bR\times M$ and $\theta$ is smooth on $W$, so it is sufficient to show that $W=\D$.

    Suppose not, so there is a $(\tau,p_0)\in\D-W$.
    We assume that $\tau>0$, the proof for negative $\tau$ is similar.
    Define $t_0=\inf\set{t\in\bR}[(t,p_0)\notin W]$.
    Taking a coordinate domain around $p_0$, we see that by uniqueness of ODE solutions, $\theta$ is smooth in a product neighborhood around $(0,p_0)$, so $t_0>0$.
    Since $t_0\leq\tau$ and $\D^{(p_0)}$ is an open interval containing $0$, we have that $t_0\in\D^{(p_0)}$.
    Let $q_0=\theta^{(p_0)}(t_0)$, and then again by ODE uniqueness, there is an $\epsilon>0$ and a neighborhood $U_0$ of $q_0$ such that $(-\epsilon,\epsilon)\times U_0\subseteq W$.
    We will show that using the group action, we can extend $\theta$ smoothly to a neighborhood of $(t_0,p_0)$.
    By maximality of $\theta$, this implies $(t_0,p_0)\in W$, but $W$ is open and thus this is a contradiction (since then a sequence of points outside of $W$ converge to a point in $W$).

    Choose $t_1<t_0$ such that $t_1+\epsilon>t_0$ and $\theta^{(p_0)}(t_1)\in U_0$ (which exists because $\theta^{(p_0)}$ is continuous).
    Since $t_1<t_0$, we have that $(t_1,p_0)\in W$ and so there is a product neighborhood $(t_1-\delta,t_1+\delta)\times U_1\subseteq W$.
    This implies that $\theta$ is defined and smooth on $[0,t_1+\delta)\times U_1$, by definition of $W$.
    Because $\theta(t_1,p_0)\in U_0$, we can choose $U_1$ small enough so that $\theta$ maps $\set{t_1}\times U_1$ into $U_0$.
    Define $\tilde\theta\colon[0,t_1+\epsilon)\times U_1\to M$ by
    $$ \tilde\theta(t,p) = \cases{\theta_t(p) & $0\leq t<t_1$\cr \theta_{t-t_1}\circ\theta_{t_1}(p) & $t_1-\epsilon<t<t_1+\epsilon$} . $$
    The group action laws show that the two definitions agree where they overlap, and our choices of $U_1,t_1,\epsilon$ ensure that this is a smooth map.

    Thus $t\mapsto\tilde\theta(t,p)$ is an integral curve of $V$ for each $p\in U_1$, so $\tilde\theta$ is a smooth extension of $\theta$ to a neighborhood of $(t_0,p_0)$.
    This is a contradiction, so $W=\D$.
    Thus $\theta$ is the unique maximal flow with infinitesimal generator $V$.

    We now show that $M_t$ is open.
    Recall that
    $$ M_t = \set{p\in M}[(t,p)\in\D] , $$
    so clearly $M_t$ is open since $\D$ is (it is the preimage of $\D$ under the map $p\mapsto(t,p)$).
    We also have
    $$ p\in M_t \implies t\in\D^{(p)} \implies \D^{(\theta_t(p))} = \D^{(p)} - t \implies -t \in \D^{(\theta_t(p))} \implies \theta_t(p) \in M_{-t} . $$
    Thus $\theta_t$ is a map $M_t\to M_{-t}$, and smooth since $\theta$ is.
    The group laws show $\theta_t\circ\theta_{-t}$ is the identity, so that $\theta_t$ is then a diffeomorphism $M_t\to M_{-t}$ as desired.
    \qed

\eproof

\bdefn

    The unique maximal flow with infinitesimal generator $V$, whose existence and uniqueness was proven above, is called the {\emphcolor flow generated by $V$},
    or just the {\emphcolor flow of $V$}.

\edefn

\bprop

    Suppose $F\colon M\to N$ is smooth, $X\in\X(M)$ and $Y\in\X(N)$.
    Let $\theta$ be the flow of $X$ and $\eta$ the flow of $Y$.
    If $X$ and $Y$ are $F$-related then for each $t\in\bR$, $F(M_t)\subseteq N_t$ and $\eta_t\circ F=F\circ\theta_t$ on $M_t$:

    \commsq{M_t&N_t\cr M_{-t}&N_{-t}}
        {F}{F}{\theta_t}{\eta_t}

\eprop

\bproof

    For any $p\in M$, $F\circ\theta^{(p)}$ is an integral curve of $Y$ starting at $F(p)$.
    By uniqueness of integral curves, the maximal integral curve $\eta^{(F(p))}$ must be defined at least on $\D^{(p)}$ and $F\circ\theta^{(p)}=\eta^{(F(p)}$
    on $\D^{(p)}$.
    This means that
    $$ p\in M_t \implies t \in \D^{(p)} \implies t\in\D^{(F(p))} \implies F(p) \in N_t , $$
    so $F(M_t)\subseteq N_t$.
    And also
    $$ F\circ\theta^{(p)}(t) = \eta^{(F(p))}(t) \implies F\circ\theta_t(p) = \eta_t(F(p)) = \eta_t\circ F(p) $$
    for all $p\in M_t$.
    \qed

\eproof

\bcoro

    Let $F\colon M\to N$ be a diffeomorphism.
    If $X\in\X(M)$ has flow $\theta$, then the flow of $F_*X$ is $\eta_t=F\circ\theta_t\circ F^{-1}$ with domain $N_t=F(M_t)$.

\ecoro

\bproof

    By the above proposition, and going in reverse as $F_*X$ and $X$ are $F^{-1}$-related.
    \qed

\eproof

\bdefn

    A vector field is {\emphcolor complete} if its flow is global.
    Equivalently if for every $p\in M$, it has an integral curve starting at $p$ defined on all of $\bR$.

\edefn

\blemm[title=Uniform time lemma]

    Let $V\in\X(M)$ with flow $\theta$.
    Suppose there is an $\epsilon>0$ such that for every $p\in M$, the domain of $\theta^{(p)}$ contains $(-\epsilon,\epsilon)$.
    Then $V$ is complete.

\elemm

\bproof

    Suppose $V$ is not complete, so for some $p\in M$ we have that $\D^{(p)}$ is not all of $\bR$.
    Without loss of generality, we assume that $\D^{(p)}$ is bounded above, let $b$ be its supremum.
    Then let $t_0\in\D^{(p)}$ such that $b-\epsilon<t_0<b$, and let $q\in\theta^{(p)}(t_0)$.
    They hypothesis implies that $\theta^{(q)}(t)$ is defined at least for $t\in(-\epsilon,\epsilon)$, so define $\gamma\colon(-\epsilon,t+\epsilon)\to M$ by
    $$ \gamma(t) = \cases{\theta^{(p)}(t) & $-\epsilon<t<b$\cr \theta^{(q)}(t-t_0) & $t_0-\epsilon<t<t_0+\epsilon$} . $$
    By the group action laws, the two definitions agree on their overlap: $\theta^{(q)}(t-t_0)=\theta(t-t_0,\theta(t_0,p))=\theta(t,p)$.
    By the translation law, $\gamma$ is an integral curve for $V$ starting at $p$.
    But since $t_0+\epsilon>b$, this is a contradiction.
    \qed

\eproof

\bthrm

    Every compactly supported smooth vector field is complete.

\ethrm

\bproof

    Let $V$ be compactly supported on $M$, and let $K=\supp V$ be its support.
    For $p\in K$ there is a neighborhood $U_p$ of $p$ and a positive number $\epsilon_p$ such that the flow of $V$ is defined on at least $(-\epsilon_p,\epsilon_p)\times U_p$,
    since the flow domain is open.
    By compactness, finitely many $U_p$ cover $K$: say $U_{p_1},\dots,U_{p_n}$ cover $K$.
    Let $\epsilon$ be the minimum of $\epsilon_{p_1},\dots,\epsilon_{p_n}$, so that the flow is defined on $(-\epsilon,\epsilon)\times K$.
    Since $V$ is zero outside of $K$, every integral curve starting outside of $K$ is constant and thus can be defined on all of $\bR$.
    Thus by the above lemma, $V$ is complete.
    \qed

\eproof

\bcoro

    Every vector field defined on a compact manifold is complete.

\ecoro

\bthrm

    Every left-invariant vector field on a Lie group is complete.

\ethrm

\bproof

    Let $G$ be a Lie group and $X\in\lie(G)$ with flow $\theta$.
    There is some $\epsilon>0$ such that $\theta^{(e)}$ is defined on $(-\epsilon,\epsilon)$.
    Let $g\in G$ so that $X$ is $L_g$-related to itself, so $L_g\circ\theta^{(e)}$ is an integral curve of $g$ defined on $(-\epsilon,\epsilon)$.
    Thus $\theta^{(g)}$ is defined on at least $(-\epsilon,\epsilon)$ for all $g\in G$, thus $X$ is complete by the uniform time lemma.
    \qed

\eproof

\subsection{Flowouts}

Recall that for an immersed submanifold $S\subseteq M$, $X\in\X(M)$ is {\emphcolor tangent to $S$} if $X_p\in T_pS\subseteq T_pM$ for all $p\in S$.
Equivalently, iff it is $\iota$-related to a smooth vector field on $U$ denoted $X|_S\in\X(S)$.

\bdefn

    For an immersed submanifold $S\subseteq M$ and $X\in\X(M)$, we say that $X$ is {\emphcolor nowhere tangent to $S$} if
    $X_p\notin T_pS$ for all $p\in S$.

\edefn

\bthrm[title=Flowout theorem]

    Let $M$ be a smooth manifold, $S\subseteq M$ be an embedded $k$-dimensional submanifold, and $V\in\X(M)$ nowhere tangent to $S$.
    Let $\theta\colon\D\to M$ be $V$'s flow, let $\O=(\bR\times S)\cap\D$, and $\Phi=\theta|_\O$.
    Then:
    \benum
        \item $\Phi\colon\O\to M$ is an immersion,
        \item $\ipdv t\in\X(O)$ is $\Phi$-related to $V$,
        \item There exists a smooth positive function $\delta\colon S\to\bR$ such that the restriction of $\Phi$ to $\O_\delta$ is injective, where
        $$ \O_\delta = \set{(t,p)\in\O}[\abs t<\delta(p)] \subseteq \O $$
        is a flow domain.
        Thus $\Phi(\O_\delta)$ is an immersed submanifold of $M$ containing $S$, and $V$ is tangent to this submanifold.
        \item If $S$ has codimension $1$, then $\Phi|_{\O_\delta}$ is a diffeomorphism onto an open submanifold of $M$.
    \eenum

\ethrm

\bnote

    The submanifold $\Phi(\O_\delta)\subseteq M$ is called a {\emphcolor flowout from $S$ along $V$}.

\enote

\bproof

    First we prove (2).
    Fix $p\in S$ and let $\sigma\colon\D^{(p)}\to\bR\times S$ be the curve $\sigma(t)=(t,p)$.
    Then $\Phi\circ\sigma(t)=\theta(t,p)$ is an integral curve of $V$, so for any $t_0\in\D^{(p)}$ we have that
    $$ d\Phi_{(t_0,p)}\parens{\evpdv t_{(t_0,p)}} = (\Phi\circ\sigma)'(t_0) = V_{(t_0,p)} , $$
    as desired.

    Now we prove (1).
    The restriction of $\Phi$ to $\set0\times S$ is the restriction of the diffeomorphism $\set0\times S\cong S$ with the embedding $S\inj M$
    (since it is the map $\Phi(0,p)=\theta(0,p)=p$), and thus an embedding.
    So the restriction of $d\Phi_{(0,p)}$ to $T_pS$ (viewed as a subspace of $T_{(0,p)}\O\cong T_0\bR\oplus T_pS$) is the inclusion $T_pS\inj T_pM$.
    Then $d\Phi_{(0,p)}$ maps the basis of $T_{(0,p)}\O$, $(\ievpdv t_{(0,p)},E_1,\dots,E_k)$ with $E_1,\dots,E_k\in T_pS$, to $(V_p,E_1,\dots,E_k)$.
    Since $V$ is nowhere tangent to $S$, $V_p$ is not in the span of $E_1,\dots,E_k$ and thus this remains linearly independent.
    Thus $d\Phi_{(0,p)}$ is an injection.

    To show that $d\Phi$ is injective elsewhere, let $(t_0,s_0)\in\O$ and let $\tau_{t_0}\colon\O\to\bR\times S$ be the translation $\tau_{t_0}(t,p)=(t+t_0,p)$.
    By the group law for $\theta$, the following diagram commutes (the horizontal maps may only be defined in a neighborhood of $(0,p_0)$ and $p_0$ respectively):

    \commsq{\O&\O\cr M&M}
        {\tau_{t_0}}{\theta_{t_0}}{\Phi}{\Phi}

    Indeed, $\theta_{t_0}\circ\Phi(t,p)=\theta_{t+t_0}(p)=\Phi\circ\tau_{t_0}(t,p)$.
    The horizontal maps are local diffeomorphisms ($\tau_{t_0}$ clearly is, $\theta_t$ is by the flow theorem), and thus the induced commutative square

    \commsq{T_{(0,p_0)}\O&T_{(t_0,p_0)}\O\cr T_{p_0}M&T_{\Phi(t_0,p_0)}M}
        {d(\tau_{t_0})_{(0,p_0)}}{d(\theta_{t_0})_{p_0}}{d\Phi_{(0,p_0)}}{d\Phi_{(t_0,p_0)}}

    has invertible horizontal maps.
    Thus the two vertical maps have the same rank, and so $\Phi$ has constant rank.
    Since it has full rank at $(0,p_0)$, it has full rank everywhere and is thus an immersion.

    Now we prove (3).
    Let $p_0\in S$ and let $(U,(x^i))$ be a chart for $S$ centered at $p_0$, so that $U\cap S$ is the set $x^{k+1}=\cdots=x^n=0$.
    Because $V$ is not tangent to $S$, one of the $n-k$ last components of $V_{p_0}$ is nonzero, say $V^j_{p_0}$.
    Shrinking $U$ if necessary, we may assume that $\abs{V^j(p)}\geq c$ for all $p\in U$.

    Now $\Phi^{-1}U$ is open in $\bR\times S$, so we may choose $\epsilon_{p_0}>0$ and a neighborhood $W_{p_0}$ of $p_0$ such that
    $(-\epsilon_{p_0},\epsilon_{p_0})\times W_{p_0}\subseteq\O$ and $\Phi((-\epsilon_{p_0},\epsilon_{p_0})\times W_{p_0})\subseteq U$.
    Let us write $\Phi$ in these local coordinates as
    $$ \Phi(t,p) = (\Phi^1(t,p),\dots,\Phi^n(t,p)) . $$
    Because $\Phi$ is the restriction of a flow, we have that $\Phi'(t,p)=V(\Phi(t,p))$, so in particular
    $$ \pdvof{\Phi^j}t(t,p) = V^j(\Phi(t,p)), $$
    and thus by the fundamental theorem of calculus, $\abs{\Phi^j(t,p)}\geq c\abs t$.
    So for $(t,p)\in(-\epsilon_{p_0},\epsilon_{p_0})\times W_{p_0}$ we have that $\Phi(t,p)\in S$ iff $t=0$ (otherwise its $j$th component is nonzero).

    Now let $\set{\psi_p}_{p\in S}$ be a partition of unity subordinate to $\set{W_p}_{p\in S}$ and define $f\colon S\to\bR$ by
    $$ f(q) = \sum_{p\in S}\epsilon_p\psi_p(q) . $$
    This is smooth and positive.
    For $q\in S$ there are finitely many $p$ such that $\psi_p(q)>0$, and if we let $\epsilon_{p_0}$ is maximal among those $p$s then
    $$ f(q) \leq \epsilon_{p_0}\sum_{p\in S}\psi_p(q) = \epsilon_{p_0} . $$
    Thus if $(t,q)\in\O$ and $\abs t<f(q)$ then $(t,q)\in(-\epsilon_{p_0},\epsilon_{p_0})\times W_{p_0}$ so that $\Phi(t,q)\in S$ iff $t=0$.

    So define $\delta=\frac12f$, and $\O_\delta$ as above:
    $$ \O_\delta = \set{(t,p)\in\O}[\abs t<\delta(p)] . $$
    We show that $\Phi$ is injective on $\O_\delta$ (and thus its image is an immersed submanifold).
    Suppose $\Phi(t,q)=\Phi(t',q')$, and we can assume $f(q')\leq f(q)$.
    Our assumption means that $\theta_t(q)=\theta_{t'}(q')$, and thus $\theta_{t-t'}(q)=q'\in S$ by the group laws.
    Now we have
    $$ \abs{t-t'} \leq \abs t + \abs{t'} \leq \frac12{f(q)+f(q')} \leq f(q) , $$
    and thus as we showed this means $\theta_{t-t'}(q)=\Phi(t-t',q)\in S$ iff $t=t'$.
    It then follows that $q=q'$ (because we can recover the initial point from an integral flow).

    Finally for (4), if $S$ has codimension $1$ then $\Phi|_{\O_\delta}$ is an injective smooth immersion between manifolds of the same dimension, and thus an
    embedding, and a diffeomorphism onto an open subset (the only submanifolds of codimension zero are open).
    \qed

\eproof

\bdefn

    Let $V\in\X(M)$, a {\emphcolor singular point} of $V$ is a point $p\in M$ such that $V_p=0$.
    Otherwise $p\in M$ is a {\emphcolor regular point}.

\edefn

\bprop

    Let $V\in\X(M)$, and $\theta\colon\D\to M$ be its flow.
    If $p\in M$ is a singular point of $V$, then $\D^{(p)}=\bR$ and $\theta^{(p)}$ is the constant curve $\theta^{(p)}(t)=p$.
    If $p$ is a regular point, then $\theta^{(p)}\colon\D^{(p)}\to M$ is a smooth immersion.

\eprop

\bproof

    If $V_p=0$, then the constant curve $\gamma\colon\bR\to M$ given by $\gamma(t)=p$ is clearly integral.
    By uniqueness and maximality, it must be equal to $\theta^{(p)}$.

    Suppose that $\theta^{(p)}$ is not an immersion; we will show that $p$ is singular.
    Since it is not an immersion, ${\theta^{(p)}}'(s)=0$ for some $s\in\D^{(p)}$.
    Let $q=\theta^{(p)}(s)$ so that $V_q=0$ and $q$ is singular.
    But then $\theta^{(q)}$ is constant, and
    $$ \theta^{(p)}(t) = \theta(t,p) = \theta(t-s,\theta(s,p)) = \theta(t-s,q) = q, $$
    so that $p=\theta^{(p)}(0)=q$ and so $p$ is singular.
    \qed

\eproof

\bdefn

    If $\theta\colon\D\to M$ is a flow, a point $p\in M$ is an {\emphcolor equilibrium point} of $\theta$ if $\theta(t,p)=p$ for all $t\in\D^{(p)}$.

\edefn

The above proposition shows that if $\theta$ is $V$'s flow, then $\theta$'s equilibrium points are precisely $V$'s singular points.

\bthrm[title=Canonical form of flow]

    Let $V\in\X(M)$ and $p\in M$ a regular point of $V$.
    Then there exist smooth coordinates $(s^i)$ in a neighborhood of $p$ such that $V$ has the coordinate representation $\ipdv{s^1}$.
    If $S\subseteq M$ is an embedded hypersurface (embedded submanifold of codimension $1$) with $p\in S$ and $V_p\notin T_pS$, then the coordinates can
    be chosen so that $s^1$ is a local defining map for $S$.

\ethrm

\bnote

    If $S\subseteq M$ is embedded, a {\emphcolor local defining map} for $S$ is a smooth map $\Phi\colon U\to N$ with $U\subseteq M$ open, such that
    $S\cap U$ is a regular level set of $\Phi$.

\enote

\bproof

    If no $S$ is given, let $(U,(x^i))$ be any smooth coordinates centered at $p$ and take $S\subseteq U$ to be the hypersurface defined by $x^j=0$
    where $j$ is any index where $V^j(p)\neq0$.

    Regardless, we now just assume we have an embedded hypersurface containing $p$ where $V_p\notin T_pS$.
    We can potentially shrink $S$ so that $V$ is nowhere tangent to it.
    The flowout theorem gives us a flow domain $\O_\delta\subseteq\bR\times S$ such that the flow of $V$ restricts to a diffeomorphism from $\O_\delta$
    to an open subset $W\subseteq M$ containing $S$.
    Take a product neighborhood $(0,p)\in(-\epsilon,\epsilon)\times W$ in $\O_\delta$, and choose a smooth local parameterization $X\colon\Omega\to S$
    where $\Omega$ is an open subset of $\bR^{n-1}$ with coordinates denoted by $s^2,\dots,s^n$, and $X$'s image is contained in $W_0$.

    Then define $\Psi\colon(-\epsilon,\epsilon)\times\Omega\to M$ by
    $$ \Psi(t,s^2,\dots,s^n) = \Phi(t,X(s^2,\dots,s^n)) $$
    is a diffeomorphism onto a neighborhood of $p$ in $M$.
    Because $(t,s^2,\dots,s^n)\mapsto(t,X(s^2,\dots,s^n))$ pushes $\ipdv t$ forward to itself, and $\Phi_*(\ipdv t)=V$ (since by the flowout theorem,
    $\ipdv t$ is $\Phi$-related to $V$), it follows that $\Psi_*(\ipdv t)=V$.

    Thus $\Psi^{-1}$ is a smooth coordinate chart in which $V$ has coordinate representation $\ipdv t$.
    Rename $t$ to $s^1$ to complete the proof: $s^1=0$ gives $W_0$, so it is a local defining map for $S$.
    \qed

\eproof

\subsection{Lie derivatives}

Tangent spaces allow us to make sense of directional derivatives of real-valued functions on manifolds; an element of $T_pM$ is viewed as a directional
derivative at $p$.
But what about the directional derivative of a vector field.
In Euclidean space, if $W$ is a vector field on $\bR^n$ then we could define its directional derivative in the direction of $v\in T_p\bR^n$ to be
$$ D_vW(p) = \evdrv t_0 W_{p+tv} = \lim_{t\to0}\frac{W_{p+tv}-W_p}t . $$
It is easy to see that
$$ D_vW(p) = D_vW^i(p)\evpdv{x^i}_p , $$
where $W^i$ are the components of $W$ and $D_vW^i$ is the usual directional derivative.

In smooth manifolds, $W_{p+tv}$ makes no sense, but we could instead replace $p+tv$ with a curve $\gamma$ starting at $p$ with initial velocity $v$.
The issue is then that $W_{\gamma(t)}$ and $W_p$ lie in different tangent spaces, and thus there is no good way of combining them.

Instead we can replace $v\in T_pM$ with a vector field $V\in\X(M)$, and so we can use the flow of $V$ to push values of $W$ back to $p$ and then differentiate.
This motivates the following.

\bdefn

    Let $M$ be smooth and without boundary, $V\in\X(M)$ with flow $\theta$.
    For any other smoot vector field $W$, define a rough vector field on $M$, denoted $\L_VW$ and called the {\emphcolor Lie derivative} of $W$ with respect to
    $V$, to be
    $$ (\L_VW)_p = \evdrv t_0 d(\theta_{-t})_{\theta_t(p)}(W_{\theta_t(p)}) = \lim_{t\to0}\frac{d(\theta_{-t})_{\theta_t(p)}(W_{\theta_t(p)})-W_p}t , $$
    provided the derivative exists.
    For small $t\neq0$, $\theta_t$ is defined in a neighborhood of $p$, and $\theta_{-t}$ is its inverse, so both $d(\theta_{-t})_{\theta_t(p)}(W_{\theta_t(p)})$
    and $W_p$ are both elements of $T_pM$.

\edefn

Note that we are differentiating a function $\bR\to T_pM$ here.
Since $T_pM$ is Euclidean, this is well-defined as any linear isomorphism from $\bR^n$ to it will give the same limit.

\blemm

    The Lie derivative of $W$ with respect to $V$ exists for all $p\in M$ and $\L_VW$ forms a smooth vector field.

\elemm

\bproof

    Let $p\in M$ and $(U,(x^i))$ be a smooth chart containing $p$.
    Choose an open interval $0\in J_0$, and an open set $U_0\subseteq U$ such that $\theta$ maps $J_0\times U_0$ into $U$.
    For $(t,x)\in J_0\times U_0$, write the components of $\theta$ as $(\theta^1(t,x),\dots,\theta^n(t,x))$.
    Then for any $(t,x)\in J_0\times U_0$, the matrix of $d(\theta_{-t})_{\theta_t(x)}\colon T_{\theta_t(x)}M\to T_xM$ is
    $$ \parens{\pdvof{\theta^i}{x^j}(-t,\theta(t,x))} . $$

    Thus
    $$ d(\theta_{-t})_{\theta_t(x)}(W_{\theta_t(x)}) = \pdvof{\theta^i}{x^j}(-t,\theta(t,x))W^j(\theta(t,x))\evpdv{x^i}_x . $$
    Because $\theta^i$ and $W^j$ are smooth, the coefficient of $\ievpdv{x^i}_x$ depends smoothly on $(t,x)$.
    It then follows that $(\L_VW)_x$, which is obtained by differentiating this and setting $t=0$, exists for each $x\in U_0$ and depends smoothly
    on $x$.
    \qed

\eproof

\bexam

    Let $v\in\bR^n$ and $W$ a smooth vector field on $\bR^n$.
    Then $D_vW(p)$, which we defined above, is just $(\L_VW)_p$, where $V$ is the vector field $v^i\ievpdv{x^i}$ with constant coefficients.

\eexam

The definition of the Lie derivative is unwieldy for computations, but it turns out that it actually is just the Lie bracket in disguise.

\bthrm

    For $V,W\in\X(M)$, $\L_VW=[V,W]$.

\ethrm

\bproof

    Let $\c R(V)\subseteq M$ be the set of regular points of $V$ (points $p$ where $V_p\neq0$).
    By continuity, $\c R(V)$ is open in $M$, and its closure is $V$'s support.
    We now show that $(\L_VW)_p=[V,W]_p$ for all $p\in M$ by considering three cases.

    \benum
        \item If $p\in\c R(V)$, then we can find a local chart $(U,(u^i))$ around $p$ where $V$ has the coordinate representation $V=\pdv{u^1}$.
        In these coordinates, $V$'s flow is then $\theta_t(u)=(u^1+t,u^2,\dots,u^n)$ (this boils down to the Euclidean case discussed in above examples).
        Then for each fixed $t$, the matrix $d(\theta_{-t})_{\theta_t(x)}$ in these coordinates is the identity.
        Thus
        $$ d(\theta_{-t})_{\theta_t(u)}(W_{\theta_t(u)}) = d(\theta_{-t})_{\theta_t(u)}\parens{W^j(u^1+t,u^2,\dots,u^n)\evpdv{u^j}_{\theta_t(u)}}
        = W^j(u^1+t,u^2,\dots,u^n)\evpdv{u^j}_u . $$
        And so by the definition of the Lie derivative,
        $$ (\L_VW)_u = \evdrv t_0 W^j(u^1+t,u^2,\dots,u^n)\evpdv{u^j}_u = \pdvof{W^j}{u^1}(u^1,\dots,u^n)\evpdv{u^j}_u . $$
        Conversely,
        $$ [V,W] = \bracks{\pdv{u^1},W^j\pdv{u^j}} = \parens{V^i\pdv{W^j}{u^i} - W^i\pdv{V^j}{u^i}}\pdv{u^j} = \parens{W^j}{u^1}\pdv{u^j} , $$
        since $V^1=1$ is constant and all other components are zero.

        \item If $p\in\supp V$ then by continuity we have the desired result.

        \item If $p\notin\supp V$ then $V=0$ on a neighborhood of $p$, and so $[V,W]_p=0$.
        Then $\theta_t$ is the identity map in a neighborhood of $p$, so $d(\theta_{-t})_{\theta_t(p)}(W_{\theta_t(p)})=W_p$ and so $(\L_VW)_p=0$ as well.
        \qed
    \eenum

\eproof

Thus we can extend the notion of a Lie derivative to manifolds with boundary simply by $\L_VW=[V,W]$.
This is equivalent to embedding $M$ into a smooth manifold without boundary $\tilde M$ (e.g.\ its double), extending $V$ and $W$ smoothly,
and computing the Lie derivative there.
By the above theorem, this just gives $[V,W]$, so is independent on the extension.

This also gives a geometric interpretation of the Lie bracket as the Lie derivative, i.e.\ a directional derivative.

\bcoro

    Let $V,W,X\in\X(M)$, then
    \benum
        \item $\L_VW=-\L_WV$,
        \item $\L_V[W,X]=[\L_VW,X]+[W,\L_VX]$,
        \item $\L_{[V,W]}X=\L_V\L_WX-\L_W\L_VX$,
        \item If $g\in C^\infty(M)$, then $\L_V(gW)=(Vg)W+g\L_VW$,
        \item If $F\colon M\to N$ is a diffeomorphism, $F_*(\L_VW)=\L_{F_*V}F_*W$.
    \eenum

\ecoro

\subsection{Commuting vector fields}

\bdefn

    Two vector fields $V,W\in\X(M)$ {\emphcolor commute} if $VWf=WVf$ for all $f\in C^\infty(M)$, or equivalently if
    $[V,W]=0$.

\edefn

\bdefn

    If $\theta$ is a smooth flow and $W\in\X(M)$, then $W$ is {\emphcolor invariant under $\theta$} if $W$ is
    $\theta_t$-related to itself for all $t$.
    More precisely, $W|_{M_t}$ is $\theta_t$-related to $W_{M_{-t}}$:
    $$ d(\theta_t)_p(W_p) = W_{\theta(t,p)} $$
    for all $(t,p)$ in the flow domain.

\edefn

\bthrm

    For $V,W\in\X(M)$, the following are equivalent:
    \benum
        \item $V$ and $W$ commute,
        \item $V$ is invariant under $W$'s flow,
        \item $W$ is invariant under $V$'s flow.
    \eenum

\ethrm

\bproof

    Let $\theta$ denote $V$'s flow.
    If $W$ is invariant under it, then we see that
    $$ d(\theta_{-t})_{\theta(t,p)}(W_{\theta(t,p)}) = W_p , $$
    and by the definition of the Lie derivative $\L_VW=0$ so $[V,W]=0$.

    Now if $\L_V(W)=[V,W]=0$ then the derivative of $d(\theta_{-t}))_{\theta(t,p)}(W_{\theta(t,p)})$ is zero
    and thus it is constant over $t$.
    Plugging in $t=0$ we see that it is equal to $W_p$, and thus we have the desired result.
    \qed

\eproof

\bcoro

    Every vector field is invariant under its own flow.

\ecoro

We would like to discuss what it means for two flows to commute.
For global flows, we could just require that $\psi_s\circ\theta_t=\theta_t\circ\psi_s$, but for local flows this may not
always be defined.
But even if we require that this only hold when both sides are defined, commuting vector fields may not have commuting flows.
The issue is that if $\theta_t\circ\psi_s(p)$ is defined for $t=t_0,s=s_0$, then it must be defined for all $t$ in some
interval containing $0$ and $t_0$, but $\psi_s$ may not be defined all the way to $s=t_0$.

\bdefn

    If $\theta,\psi$ are smooth flows on $M$, then we say they {\emphcolor commute} if whenever $p\in M$ and
    $J,K$ are open intervals containing $0$ such that one of $\theta_t\circ\psi_s(p),\psi_s\circ\theta_t(p)$ exist
    for all $(t,s)\in J\times K$, then both exist and are equal.

\edefn

\bthrm

    Smooth vector fields commute iff their flows commute.

\ethrm

\bproof

    Let $V,W\in\X(M)$ and $\theta,\psi$ be their respective flows.
    Suppose $V$ and $W$ commute and $J,K$ are intervals such that $\psi_s\circ\theta_t(p)$ is defined for all
    $(t,s)\in J\times K$.
    Since $V$ is invariant under $\psi$ by assumption, let $s\in K$ and consider the curve $\gamma\colon J\to M$
    given by $\gamma(t)=\psi_s\circ\theta_t(p)$.
    Then $\gamma(0)=\psi_s(p)$ and its velocity at $t\in J$ is
    $$ \gamma'(t) = d(\psi_s)_{\theta(t,p)}({\theta^{(p)}}'(t)) = d(\psi_s)_{\theta(t,p)}(V_{\theta(t,p)})
    = V_{\psi(s,\theta(t,p))} = V_{\gamma(t)} , $$
    so that $\gamma(t)$ is integral for $V$.
    Then
    $$ \gamma(t) = \theta^{(\psi(s,p))}(t) = \theta_t(\psi_s(p)) , $$
    so that $\theta,\psi$ commute.

    Now if the flows commute, let $p\in M$ and $\epsilon>0$ small enough so that $\psi_s\circ\theta_t(p)=\theta_t\circ\psi_s(p)$
    for $\abs t,\abs s<\epsilon$.
    Now we have
    $$ W_{\theta_t(p)} = \evdrv s_{s=0}\psi^{(\theta_t(p))}(s) = \evdrv s_{s=0}\theta_t(\psi^{(p)}(s))
    = d(\theta_t)_p(W_p) . $$
    Showing that $W$ is $\theta_t$-invariant, as desired.
    \qed

\eproof

\bdefn

    A {\emphcolor commuting frame} for a smooth manifold $M$ is a local frame $(E_i)$ such that each $E_i,E_j$ commute.

\edefn

Note that coordinate frames $\ievpdv x_p$ are commuting, as shown in our discussion of Lie brackets.
Thus if a frame doesn't commut, it cannot be a coordinate frame.

\bthrm

    Let $M$ be smooth and $(V_1,\dots,V_k)$ be linearly independent commuting local vector fields, defined on an open
    subset $W$.
    Then for every $p\in W$ there exists a smooth coordinate chart $(U,(s^i))$ centered at $p$ such that $V_i=\ipdv{s^i}$
    for all $i=1,\dots,k$.

    If $S\subseteq M$ is embedded and of codimension $k$ and $p\in S$ is such that $T_pS$ is complementary to the span
    of $V_1|_p,\dots,V_k|_s$, then the coordinate chart can be chosen so that $S\cap U$ is the slice given by
    $s^1=\cdots=s^k=0$.

\ethrm

\bproof

    If no $S$ is given, we can find one via an appropriate coordinate slice.
    Let $(U,(x^i))$ be the slice chart for $S$ centered at $p$, with $U\subseteq W$ and $S\cap U$ being given by
    $x^1=\cdots=x^k=0$.
    Our assumptions show that $V_1,\dots,V_k,\ipdv{x^{k+1}},\dots,\ipdv{x^n}$ form a local frame defined on $U$.
    Since this theorem is local, we can assume $U\subseteq\bR^n$.

    Let $\theta_i$ denote the flow of $V_i$.
    Then since the vector fields commute, there is an $\epsilon>0$ such that
    $(\theta_1)_{t_1}\circ\cdots\circ(\theta_k)_{t_k}$ is defined when $\abs{t_i}<\epsilon$ and on a neighborhood $Y$ of $p$.
    Define $\Omega\subseteq\bR^{n-k}$ by
    $$ \Omega = \set{(s^{k+1},\dots,s^n)}[(0,\dots,0,s^{k+1},\dots,s^n)\in Y] . $$
    Define $\Phi\colon(-\epsilon,\epsilon)^k\times\Omega\to U$ by
    $$ \Phi(s^1,\dots,s^k,s^{k+1},\dots,s^n) = (\theta_1)_{s^1}\circ\cdots\circ(\theta_k)_{s^k}(0,\dots,0,s^{k+1},\dots,s^n). $$
    By construction $\Phi(\set 0\times\Omega)=S\cap Y$.
    Now we will show that $\ipdv{s^i}$ is $\Phi$-related to $V_i$ for $i\leq k$.
    For any $i\leq k$ and $s_0\in(-\epsilon,\epsilon)^k\times\Omega$, we have
    $$ d\Phi_{s_0}\parens{\evpdv{s^i}_{s_0}}f = \evpdv{s^i}_{s_0}f(\Phi(s^1,\dots,s^n))
    = \evpdv{s^i}_{s_0}f((\theta_1)_{s^1}\circ\cdots\circ(\theta_k)_{s^k}(0,\dots,0,s^{k+1},\dots,s^n)) . $$
    Since the flows commute, we can reorder this composition to put $\theta_i$ first.
    And $t\mapsto(\theta_i)_t(q)$ is integral for $V_i$ so this is equal to the derivative of $f$ composed with this
    integral curve, equal to $V_i|_{\Phi(s_0)}f$ as desired.

    Now we will show that $d\Phi_0$ is invertible.
    Since $d\Phi_0(\ipdv{s^i}_0)=V_i|_p$ by our above computation, and $\Phi(0,\bar s)=(0,\bar s)$ we see that
    $$ d\Phi_0\parens{\pdv{s^i}_0} = \pdv{s^i}_p,\qquad i>k . $$
    Thus $d\Phi_0$'s image of the coordinate basis is itself a basis, so it is invertible.

    So $\Phi$ is a diffeomorphism in a neighborhood of $0$, and $\phi=\Phi^{-1}$ is a smooth coordinate chart around $p$
    where $V_i$'s coordinate representation is $\ipdv{s^i}$ and takes $S$ to the slice $s^1=\cdots=s^k=0$.
    \qed

\eproof

